%%%%%%%%%%  scan idx 211  |  folio 198  %%%%%%%%%%
tie de (0.3).

Enfin, soit $E\in\Ob\Qcoh(Z)$. On a l'isomorphisme de dualité
\[
 \mathbb R\Hom_{\Mod}(E,\mathbb Rj_{1*}(\widetilde I))
 \simeq \mathbb R\Hom_{\Mod}(j_1^*E,\widetilde I),
\]
mais comme $X$ est affine, on a (II 3.5)
\[
 \mathbb R\Hom_{\Mod}(j_1^*E,\widetilde I)
 \simeq \mathbb R\Hom_{\Qcoh}(j_1^*E,I)
 \simeq \Hom(j_1^*E,\widetilde I).
\]
On en déduit
\[
 \mathbb R\Hom_{\Mod}(E,\mathbb Rj_{1*}(\widetilde I))
 \simeq \Hom(j_1^*E,\widetilde I),
\]
d'où une suite spectrale, dont la suite exacte des termes de bas degrés donne l'isomorphisme
\[
 \Ext^2_{\Mod}(E,j_{1*}\widetilde I)
 \simeq \Hom(E,R^1j_{1*}\widetilde I).
\]
Par suite, le morphisme identique de $E=R^1j_{1*}(\widetilde I)$ fournit un élément non nul de $\Ext^2_{\Mod}(E,j_{1*}\widetilde I)$, donc le foncteur $\mathrm D^b(\Qcoh(Z))\longrightarrow \mathrm D^b(Z)$ n'est pas pleinement fidèle, ce qui achève la démonstration de (0.3).



% Additions for Expose II, Appendix II
\providecommand{\cO}{\mathcal O}
\providecommand{\cE}{\mathscr E}
\providecommand{\cF}{\mathscr F}
\providecommand{\cG}{\mathscr G}
\providecommand{\cH}{\mathcal H}
\providecommand{\cL}{\mathscr L}
\providecommand{\cP}{\mathscr P}
\providecommand{\cM}{\mathscr M}
\providecommand{\cD}{\mathcal D}
\providecommand{\R}{\mathbb R}
\providecommand{\C}{\mathbb C}
\providecommand{\N}{\mathbb N}
\providecommand{\Zell}{\mathbb Z_\ell}
\providecommand{\Qell}{\mathbb Q_\ell}
\providecommand{\C}{\mathbb C}
\providecommand{\Ch}{\operatorname{Ch}}
\providecommand{\Sym}{\operatorname{Sym}}
\providecommand{\Td}{\operatorname{Td}}
\providecommand{\Diff}{\operatorname{Diff}}
\providecommand{\Homop}{\operatorname{Hom}}
\providecommand{\Spec}{\operatorname{Spec}}
\providecommand{\pr}{\operatorname{pr}}
\providecommand{\ind}{\operatorname{ind}}
\providecommand{\soft}{\operatorname{soft}}
\providecommand{\Ban}{\operatorname{Ban}}
\providecommand{\Lop}{\operatorname{L}}
\providecommand{\im}{\operatorname{Im}}
\providecommand{\Parf}{\operatorname{Parf}}
\renewcommand{\id}{\operatorname{id}}

\newpage
